\section{Counterfactual Label Generation}
\label{sec:label-generation}

The benchmark label is defined on candidate victims rather than on completed policies. Let $r_1, r_2, \ldots$ denote a request sequence, and let $t$ index a request $r_t$ that misses when the cache is full. Write $S_t^{-}$ for the cache state immediately before the eviction choice; the candidate set $C_t = S_t^{-}$ is every resident object, so $|C_t|$ equals the cache capacity (32, 64, 128, or 256 objects in the evaluated release), counted in objects rather than bytes, with upstream object sizes not separately modeled. For any candidate $c \in C_t$, consider the intervention that evicts $c$, admits the missed request into the freed slot, and then follows a fixed continuation policy $\pi$ for the next $H$ requests, resolving every subsequent miss according to $\pi$; in the evaluated release, $\pi$ is LRU continuation rather than a Belady/offline-optimal rule. The resulting finite-horizon counterfactual loss counts those subsequent misses, so every miss contributes unit cost:
\[
\begin{aligned}
L_{t,H}^{\pi}(c)
&=
\#\{\tau \in \{t+1, \ldots, t+H\} : \\
&\qquad r_{\tau} \text{ misses in the induced trajectory}\}.
\end{aligned}
\]
The release stores this quantity as \texttt{y\_loss}, and it stores the corresponding value label as \texttt{y\_value} $= -\texttt{y\_loss}$. The triggering miss at time $t$ is common to all candidates and is therefore not the decision-discriminating part of the label.

Both columns are retained in the schema, but they are not independent
information: \texttt{y\_value} is a deterministic sign-flip of
\texttt{y\_loss}, computed by the same generation step. The two exist
because downstream modeling code is written against either a
minimization convention (\texttt{y\_loss} as a cost/regret target, the
convention used throughout the rest of this paper and the more directly
interpretable of the two -- a count of extra misses) or a maximization
convention (\texttt{y\_value} as a value/utility target, convenient for
argmax-based candidate selection or value-function-style formulations).
This paper treats \texttt{y\_loss} as canonical and uses
\texttt{y\_value} only where a specific baseline is more naturally stated
as a maximization, never reporting the same statistic under both names.

This label has two consequences for the paper. First, it creates a common numerical target for regression-style benchmark tasks. Second, it induces structured decision-level objects. For a fixed decision $(t, S_t^{-}, C_t)$, the benchmark-optimal candidate set is
\[
C_t^{\star} = \arg\min_{c \in C_t} L_{t,H}^{\pi}(c),
\]
and the per-candidate regret is
\[
\Delta_{t,H}^{\pi}(c) = L_{t,H}^{\pi}(c) - \min_{c' \in C_t} L_{t,H}^{\pi}(c').
\]
These quantities define best-candidate selection targets, pairwise preferences, and tie structure within a decision.

This supervision is a generated counterfactual label family, not an upstream annotation. It is finite-horizon, not asymptotic. It is tied to a continuation policy, not a universally optimal control objective. The label is therefore decision-aligned for the benchmarked eviction choice, but it is not an offline-optimal label in the sense of solving the full future control problem. These boundaries are central to the scientific positioning of the paper.
